% =====================================================================
% Cine-JEPA — Methods
% NeurIPS 2026 submission
% Standalone .tex fragment: include via % =====================================================================
% Cine-JEPA — Methods
% NeurIPS 2026 submission
% Standalone .tex fragment: include via % =====================================================================
% Cine-JEPA — Methods
% NeurIPS 2026 submission
% Standalone .tex fragment: include via % =====================================================================
% Cine-JEPA — Methods
% NeurIPS 2026 submission
% Standalone .tex fragment: include via \input{methods} from the main file.
% Citation keys are placeholders; resolve them in the bibliography.
% =====================================================================

\section{Method}
\label{sec:method}

\textbf{Cine-JEPA} is a joint-embedding predictive architecture
\citep{lecun2022jepa,assran2023ijepa,bardes2024vjepa,assran2025vjepa2}
that learns EEG representations by predicting the latents of masked
spatiotemporal blocks of cortical activity recorded during
naturalistic movie-watching, on top of a stimulus-locked
CorrCA-projected input. The EEG-side tokenizer, $4$D positional
embedding, and transformer trunk are taken from REVE
\citep{elouahidi2025reve} unchanged, so any gain we report is
attributable to the input pipeline, the self-supervised objective,
and our masking specialization rather than to backbone differences.

% =====================================================================
\subsection{Pretraining data and preprocessing}
\label{sec:method:data}

\paragraph{Healthy Brain Network EEG corpus.}
We pretrain on the movie-watching subset of the Healthy Brain Network
EEG dataset \citep{alexander2017hbn} accessed via the EEGDash data
loader \citep{tousignant2024eegdash} on OpenNeuro. Subjects are
children and adolescents who watched short clips while EEG was
recorded on a $129$-channel EGI GSN-HydroCel system. We use a single
clip throughout, \textit{The Present} ($203.3$\,s, $24$\,fps), to
guarantee that train, validation, and test recordings are time-aligned
to the same stimulus and therefore directly comparable under shared
movie-feature labels. HBN releases are split disjointly by subject:
releases R1--R4 (\texttt{ds005505}--\texttt{ds005508}, ${\sim}\,700$
subjects in total) for pretraining and probe training, release R5
(\texttt{ds005509}, $108$ subjects) for validation and model
selection, and release R6 (\texttt{ds005510}, $108$ subjects) as a
held-out test set used only for final evaluation. Splitting at the
release level guarantees subject-level disjointness by construction.
Recordings shorter than $10$\,s, or whose movie-onset/offset
annotations disagree with the clip duration by more than $1$\,s, are
discarded.

\paragraph{Resampling and bandpass.}
Raw recordings are resampled to a common $f_s = 200$\,Hz, bandpass
filtered to $[0.5, 99.5]$\,Hz, cast to \texttt{float32}, and clipped
to $\pm 15\sigma$ after temporal normalization to limit the influence
of movement and electrode artefacts.

\paragraph{Temporal normalization.}
Each recording $X^{(s)} \in \mathbb{R}^{C \times T_s}$ is Z-scored
\emph{per recording}: $\tilde{X}^{(s)}_{c,t} = (X^{(s)}_{c,t} -
\mu^{(s)}_c) / \sigma^{(s)}_c$, removing recording-level amplitude
differences (skin conductance, electrode impedance, instrumentation
drift) while leaving subject-trait signal recoverable by downstream
probes. Corpus-level (global) Z-scoring (the REVE convention,
\citep{elouahidi2025reve}) is reported as an ablation.

\paragraph{Spatial filtering (CorrCA).}
We then apply Correlated Component Analysis
\citep{dmochowski2012corrca} as a stimulus-locked spatial filter:
filters $W \in \mathbb{R}^{129 \times K}$ that maximize inter-subject
correlation of $W^{\!\top} X^{(s)}$ are fit on the R1--R4 recordings
of \textit{The Present}, the top $K = 5$ components are retained, and
the frozen filters are applied to all downstream data. The canonical
pipeline therefore feeds $C = 5$ CorrCA channels to the encoder;
$C = 129$ Raw is reported as an ablation.

\paragraph{Sampling and stimulus--EEG alignment.}
For each iteration we sample a contiguous $8$\,s clip ($T = 2$
windows of $4$\,s) uniformly over recordings and start times.
Cortical EEG responses to visual input are delayed: the earliest
striate response to luminance/contrast modulations sits at
${\sim}80$--$100$\,ms ($P1$ / VESPA peak) \citep{lalor2006vespa,
lalor2009vespa, lalor2012vespa}, mTRF analyses of natural stimuli
use lags out to $300$--$600$\,ms \citep{crosse2016mtrf,
diliberto2015lowfreq, lalor2010natural}, and ISC of free movie
viewing peaks at $100$--$300$\,ms \citep{dmochowski2012corrca,
dmochowski2014audience, cohen2017isc}. We therefore align labels
with a fixed visual-processing delay $\Delta = 300$\,ms: the
movie-feature label for an EEG window starting at sample $s$ is
read from frame
$k(s) = \lfloor (s / f_s - \Delta) \cdot \mathrm{fps} \rfloor$,
with boundary clamping at the movie endpoints. The exact value was
validated by a delay sweep on the training split
(App.~\ref{app:delay-sweep}).

% =====================================================================
\subsection{Cine-JEPA architecture and pretraining objective}
\label{sec:method:model}

\paragraph{Backbone instantiation.}
A clip is a tensor $X \in \mathbb{R}^{T \times C \times W}$ with
$T = 2$ non-overlapping $4$\,s windows ($W = 800$ samples at
$200$\,Hz) over $C$ channels — canonically $C = 5$ CorrCA components,
with $C = 129$ raw electrodes as an ablation. Each window is unfolded
into $P = 26$ overlapping temporal patches per channel (patch length
$\ell_p = 50$, overlap $\ell_o = 20$, stride $30$) and each patch is
linearly embedded to a $D = 64$ token, yielding $N = C \cdot T \cdot
P = 260$ tokens per clip in the canonical configuration. The context
and target encoders $f_\theta, f_{\bar\theta}$ share REVE's
transformer trunk with $L = 2$ blocks, $H = 4$ heads of dimension
$d_h = 16$, hidden width $D = 64$, and feed-forward expansion ratio
$2.66$; we refer the reader to \citep{elouahidi2025reve} for the full
block-level specification.

\paragraph{Channel-identity handling for CorrCA components.}
REVE's $4$D positional encoding takes a $3$D scalp coordinate
$(x, y, z)$ for each spatial channel together with the patch's global
time index $\tau$. For the Raw variant, $(x_c, y_c, z_c)$ is the
canonical GSN-HydroCel-$129$ coordinate of electrode $c$. The
canonical CorrCA components have no native position, so we anchor
each component $k$ to its peak-weight electrode,
\begin{equation}
  c^\star_k = \arg\max_{c} |W_{c, k}|,
  \qquad
  (x_k, y_k, z_k) = (x_{c^\star_k}, y_{c^\star_k}, z_{c^\star_k}).
\end{equation}
The CorrCA filters are dominated by a handful of central /
parieto-occipital electrodes, so the peak-weight anchor is a faithful
representative of each component's spatial support; it is also
parameter-free, montage-aware, and deterministic across runs.

\paragraph{Multi-block masking.}
We adopt V-JEPA-style multi-block masking \citep{bardes2024vjepa}
specialized to the EEG token grid. For each batch we draw $B = 4$ mask
blocks: $B_s = 2$ \emph{short-range} blocks and $B_\ell = 2$
\emph{long-range} blocks, where each block is a contiguous rectangle
on the $[C, P]$ grid. Block dimensions are sampled uniformly:
short-range blocks span $[8\%, 15\%]$ of channels and $[30\%, 60\%]$
of patches; long-range blocks span $[15\%, 35\%]$ of channels and
$[50\%, 100\%]$ of patches, jointly masking on average roughly half of
the $C \cdot P$ cells. \emph{Crucially, every $(c, p)$ cell that is
masked is masked identically across all $T$ windows}: this prevents
the predictor from exploiting the high temporal autocorrelation of
EEG to trivially copy the same channel/patch from a neighbouring
window. We reject mask draws that leave fewer than $15\%$ of tokens
visible. The context set $\mathcal{C}$ is the union of unmasked
positions; each of the $B$ blocks defines a separate prediction target
$\mathcal{M}_b$, disjoint from $\mathcal{C}$.

\paragraph{Predictor with mask tokens.}
The predictor $g_\phi$ takes the context tokens
$\{f_\theta(X)_i\}_{i \in \mathcal{C}}$, adds their position
embeddings, appends a learned mask token $\mathbf{m} \in
\mathbb{R}^{D'}$ summed with each target position embedding
$\{\mathbf{p}_j\}_{j \in \mathcal{M}_b}$, runs a transformer of the
same depth as the encoder ($L_g = L = 2$), and reads out the
final-layer states at the mask-token slots as predictions
$\hat{\mathbf{z}}_j \in \mathbb{R}^{D}$. Following I-JEPA / V-JEPA
\citep{assran2023ijepa,bardes2024vjepa}, $g_\phi$ operates at a
narrow bottleneck $D' = 24 < D = 64$ via linear input/position/output
projections, which forces the encoder to produce richer
representations rather than offloading work to the predictor.

\paragraph{Pretraining loss.}
For each masked block we minimize a smooth-L1 (Huber) regression to
the \emph{stop-gradient} target representations,
\begin{equation}
  \mathcal{L}_{\text{pred}}(X)
   = \frac{1}{|\mathcal{M}|}
     \sum_{b=1}^{B} \sum_{j \in \mathcal{M}_b}
     \mathrm{Huber}\!\Big(\hat{\mathbf{z}}_j,\;
                          \mathrm{sg}\big[f_{\bar\theta}(X)_j\big]\Big),
  \label{eq:pred-loss}
\end{equation}
with $|\mathcal{M}| = \sum_b |\mathcal{M}_b|$.

\paragraph{Anti-collapse regularization.}
A trivial constant solution drives Eq.~\eqref{eq:pred-loss} to zero,
so two complementary mechanisms prevent representational collapse.
\emph{(i) EMA target encoder.} The parameters of $f_{\bar\theta}$ are
updated as $\bar\theta \leftarrow \mu_t \, \bar\theta + (1 - \mu_t)\,
\theta$ with a cosine momentum schedule $\mu_t : 0.996 \to 1.0$ over
the full training run, and gradients flow only through $f_\theta$ and
$g_\phi$. \emph{(ii) VICReg regularization.} We additionally apply
the variance and covariance terms of VICReg \citep{bardes2022vicreg}
to the context-encoder outputs passed through a small projector
$h_\psi$,
\begin{equation}
  \mathcal{L}_{\text{vc}}(\mathbf{Z})
   = \lambda_{\sigma}\, \frac{1}{D_h}\sum_{d=1}^{D_h}
     \max\!\big(0,\, 1 - \sigma_d(\mathbf{Z})\big)
   + \lambda_{c}\, \frac{1}{D_h}\sum_{d \neq d'}
     \mathrm{Cov}_{d d'}(\mathbf{Z})^2,
   \qquad
   \mathbf{Z} = h_\psi\!\big(f_\theta(X)_{\mathcal{C}}\big),
\end{equation}
where $\sigma_d$ and $\mathrm{Cov}$ are computed across the batch and
token axes within $\mathcal{C}$ and $D_h$ is the projector output
dimension. The invariance term of VICReg is supplied implicitly by
the JEPA prediction loss in Eq.~\eqref{eq:pred-loss} and is therefore
omitted. We set $\lambda_\sigma = \lambda_c = 0.25$ and minimize
$\mathcal{L} = \mathcal{L}_{\text{pred}} + \mathcal{L}_{\text{vc}}$.

\paragraph{Optimization.}
We train for $100$ epochs with batch size $64$ on a single GPU using
AdamW \citep{loshchilov2019adamw} with peak learning rate
$5 \times 10^{-4}$, $5$-epoch linear warm-up, and cosine decay to
$10^{-6}$. Gradients on $f_\theta$, $g_\phi$, and $h_\psi$ are
clipped to global $\ell_2$ norm $1.0$; $f_{\bar\theta}$ is updated by
EMA only. At evaluation time the predictor, target encoder, and
VICReg projector are discarded. We train five seeds
$\{42, 123, 456, 789, 2025\}$ and report inter-seed variability as
our primary measure of encoder-training variance.

% =====================================================================
\subsection{Evaluation protocol}
\label{sec:method:eval}

Every method, including all baselines, is evaluated under a single
shared protocol: a deterministic clip-to-feature map followed by
closed-form linear probes and recording-level bootstrap CIs. Only
the clip-to-feature map varies between methods.

\paragraph{Per-clip representation.}
Once pretraining is complete, the context encoder $f_\theta$ is
frozen. Given an $8$\,s clip $X$ in canonical CorrCA form ($C{=}5$,
$T{=}2$, $W{=}800$), we obtain token outputs $f_\theta(X) \in
\mathbb{R}^{N \times D}$, reshape to the structured grid $[C, T, P,
D]$, and form a per-clip representation by mean-pooling over the
patch and window axes while \emph{keeping the channel axis}:
\begin{equation}
  \bar{\mathbf{h}}(X) = \mathrm{flatten}_c\!\Big(
      \tfrac{1}{T \cdot P}\!\sum_{t, p} f_\theta(X)_{c, t, p, :}
    \Big) \;\in\; \mathbb{R}^{C \cdot D}
  \;=\; \mathbb{R}^{320}.
\end{equation}
Keeping the channel axis preserves the spatial structure that CorrCA
already isolated, and makes the probe input dimension match the
linear baselines (\S\,below) for an apples-to-apples capacity
comparison.

\paragraph{Probe targets.}
We probe a deliberately broad hierarchy of targets, from low-level
sensory features through narrative structure to subject-level traits,
in order to characterise \emph{what} the representations encode rather
than only \emph{how well} they decode any one variable.
\begin{enumerate}
    \item \textit{Low-level visual features.} Per video frame we
          extract \texttt{luminance\_mean} (mean grayscale intensity)
          and \texttt{contrast\_rms} (RMS contrast), and average them
          within each window.
    \item \textit{Narrative event score} \texttt{narr} — a saliency
          score for narrative events extracted by the
          movie-annotation pipeline (\S\ref{sec:method:data}).
    \item \textit{Temporal position in the movie}
          \texttt{position\_in\_movie}, the normalised time
          $\tau / T_{\text{clip}} \in [0, 1]$, probing whether
          representations carry coarse temporal context independent
          of local stimulus content.
    \item \textit{Movie identity.} A $20$-bin top-$1$ / top-$5$
          classification over $20$ equal-width position bins, used as
          a coarser temporal probe.
    \item \textit{Subject traits.} \texttt{age} as a continuous
          regression target and biological \texttt{sex} as a binary
          classification target.
\end{enumerate}
For all probes, continuous targets are predicted by ridge regression
$\mathrm{Ridge}(\alpha{=}1)$ in closed form, binary targets by
$\ell_2$-regularized logistic regression $\mathrm{LogReg}(C{=}1)$
fitted with L-BFGS, and the $20$-bin movie-identity task by
multinomial $\ell_2$-regularized logistic regression. Regression
metrics are Pearson $r$ and (where reported) $R^2$; classification
metrics are balanced accuracy and ROC-AUC; movie-identity is reported
as top-$1$ and top-$5$ accuracy. All standardization, label
medians/means, PCA bases, and CorrCA filters are fit on training
recordings only.

\paragraph{Sampling and clip budget.}
Each recording is read $n_{\text{pass}} = 20$ times with independent
random clip starts, yielding $14{,}000$ train, $2{,}160$ val, and
$2{,}160$ test clips per encoder; val clips are used only for
protocol audits.

\paragraph{Baselines.}
Under the same protocol we compare to a linear TRF
\citep{crosse2016mtrf}, fully-supervised end-to-end EEGNet
\citep{lawhern2018eegnet}, BIOT \citep{yang2023biot}, and REVE
\citep{elouahidi2025reve}, and linear-probe + full-finetune
evaluations of LaBraM \citep{jiang2024labram}, NeuroLM
\citep{jiang2024neurolm}, and EEGPT \citep{wang2024eegpt}, all under
the same $8$\,s clip definition and R6 test split.

\paragraph{Implementation details.}
Cine-JEPA is implemented in PyTorch~$2$ \citep{paszke2019pytorch} with
\texttt{einops} for tensor reshaping, and we train each run on a
single NVIDIA GPU. All probes use scikit-learn
\citep{pedregosa2011sklearn}. Code, pretrained checkpoints, and the
movie-annotation pipeline used to compute the probe targets will be
released upon publication.
 from the main file.
% Citation keys are placeholders; resolve them in the bibliography.
% =====================================================================

\section{Method}
\label{sec:method}

\textbf{Cine-JEPA} is a joint-embedding predictive architecture
\citep{lecun2022jepa,assran2023ijepa,bardes2024vjepa,assran2025vjepa2}
that learns EEG representations by predicting the latents of masked
spatiotemporal blocks of cortical activity recorded during
naturalistic movie-watching, on top of a stimulus-locked
CorrCA-projected input. The EEG-side tokenizer, $4$D positional
embedding, and transformer trunk are taken from REVE
\citep{elouahidi2025reve} unchanged, so any gain we report is
attributable to the input pipeline, the self-supervised objective,
and our masking specialization rather than to backbone differences.

% =====================================================================
\subsection{Pretraining data and preprocessing}
\label{sec:method:data}

\paragraph{Healthy Brain Network EEG corpus.}
We pretrain on the movie-watching subset of the Healthy Brain Network
EEG dataset \citep{alexander2017hbn} accessed via the EEGDash data
loader \citep{tousignant2024eegdash} on OpenNeuro. Subjects are
children and adolescents who watched short clips while EEG was
recorded on a $129$-channel EGI GSN-HydroCel system. We use a single
clip throughout, \textit{The Present} ($203.3$\,s, $24$\,fps), to
guarantee that train, validation, and test recordings are time-aligned
to the same stimulus and therefore directly comparable under shared
movie-feature labels. HBN releases are split disjointly by subject:
releases R1--R4 (\texttt{ds005505}--\texttt{ds005508}, ${\sim}\,700$
subjects in total) for pretraining and probe training, release R5
(\texttt{ds005509}, $108$ subjects) for validation and model
selection, and release R6 (\texttt{ds005510}, $108$ subjects) as a
held-out test set used only for final evaluation. Splitting at the
release level guarantees subject-level disjointness by construction.
Recordings shorter than $10$\,s, or whose movie-onset/offset
annotations disagree with the clip duration by more than $1$\,s, are
discarded.

\paragraph{Resampling and bandpass.}
Raw recordings are resampled to a common $f_s = 200$\,Hz, bandpass
filtered to $[0.5, 99.5]$\,Hz, cast to \texttt{float32}, and clipped
to $\pm 15\sigma$ after temporal normalization to limit the influence
of movement and electrode artefacts.

\paragraph{Temporal normalization.}
Each recording $X^{(s)} \in \mathbb{R}^{C \times T_s}$ is Z-scored
\emph{per recording}: $\tilde{X}^{(s)}_{c,t} = (X^{(s)}_{c,t} -
\mu^{(s)}_c) / \sigma^{(s)}_c$, removing recording-level amplitude
differences (skin conductance, electrode impedance, instrumentation
drift) while leaving subject-trait signal recoverable by downstream
probes. Corpus-level (global) Z-scoring (the REVE convention,
\citep{elouahidi2025reve}) is reported as an ablation.

\paragraph{Spatial filtering (CorrCA).}
We then apply Correlated Component Analysis
\citep{dmochowski2012corrca} as a stimulus-locked spatial filter:
filters $W \in \mathbb{R}^{129 \times K}$ that maximize inter-subject
correlation of $W^{\!\top} X^{(s)}$ are fit on the R1--R4 recordings
of \textit{The Present}, the top $K = 5$ components are retained, and
the frozen filters are applied to all downstream data. The canonical
pipeline therefore feeds $C = 5$ CorrCA channels to the encoder;
$C = 129$ Raw is reported as an ablation.

\paragraph{Sampling and stimulus--EEG alignment.}
For each iteration we sample a contiguous $8$\,s clip ($T = 2$
windows of $4$\,s) uniformly over recordings and start times.
Cortical EEG responses to visual input are delayed: the earliest
striate response to luminance/contrast modulations sits at
${\sim}80$--$100$\,ms ($P1$ / VESPA peak) \citep{lalor2006vespa,
lalor2009vespa, lalor2012vespa}, mTRF analyses of natural stimuli
use lags out to $300$--$600$\,ms \citep{crosse2016mtrf,
diliberto2015lowfreq, lalor2010natural}, and ISC of free movie
viewing peaks at $100$--$300$\,ms \citep{dmochowski2012corrca,
dmochowski2014audience, cohen2017isc}. We therefore align labels
with a fixed visual-processing delay $\Delta = 300$\,ms: the
movie-feature label for an EEG window starting at sample $s$ is
read from frame
$k(s) = \lfloor (s / f_s - \Delta) \cdot \mathrm{fps} \rfloor$,
with boundary clamping at the movie endpoints. The exact value was
validated by a delay sweep on the training split
(App.~\ref{app:delay-sweep}).

% =====================================================================
\subsection{Cine-JEPA architecture and pretraining objective}
\label{sec:method:model}

\paragraph{Backbone instantiation.}
A clip is a tensor $X \in \mathbb{R}^{T \times C \times W}$ with
$T = 2$ non-overlapping $4$\,s windows ($W = 800$ samples at
$200$\,Hz) over $C$ channels — canonically $C = 5$ CorrCA components,
with $C = 129$ raw electrodes as an ablation. Each window is unfolded
into $P = 26$ overlapping temporal patches per channel (patch length
$\ell_p = 50$, overlap $\ell_o = 20$, stride $30$) and each patch is
linearly embedded to a $D = 64$ token, yielding $N = C \cdot T \cdot
P = 260$ tokens per clip in the canonical configuration. The context
and target encoders $f_\theta, f_{\bar\theta}$ share REVE's
transformer trunk with $L = 2$ blocks, $H = 4$ heads of dimension
$d_h = 16$, hidden width $D = 64$, and feed-forward expansion ratio
$2.66$; we refer the reader to \citep{elouahidi2025reve} for the full
block-level specification.

\paragraph{Channel-identity handling for CorrCA components.}
REVE's $4$D positional encoding takes a $3$D scalp coordinate
$(x, y, z)$ for each spatial channel together with the patch's global
time index $\tau$. For the Raw variant, $(x_c, y_c, z_c)$ is the
canonical GSN-HydroCel-$129$ coordinate of electrode $c$. The
canonical CorrCA components have no native position, so we anchor
each component $k$ to its peak-weight electrode,
\begin{equation}
  c^\star_k = \arg\max_{c} |W_{c, k}|,
  \qquad
  (x_k, y_k, z_k) = (x_{c^\star_k}, y_{c^\star_k}, z_{c^\star_k}).
\end{equation}
The CorrCA filters are dominated by a handful of central /
parieto-occipital electrodes, so the peak-weight anchor is a faithful
representative of each component's spatial support; it is also
parameter-free, montage-aware, and deterministic across runs.

\paragraph{Multi-block masking.}
We adopt V-JEPA-style multi-block masking \citep{bardes2024vjepa}
specialized to the EEG token grid. For each batch we draw $B = 4$ mask
blocks: $B_s = 2$ \emph{short-range} blocks and $B_\ell = 2$
\emph{long-range} blocks, where each block is a contiguous rectangle
on the $[C, P]$ grid. Block dimensions are sampled uniformly:
short-range blocks span $[8\%, 15\%]$ of channels and $[30\%, 60\%]$
of patches; long-range blocks span $[15\%, 35\%]$ of channels and
$[50\%, 100\%]$ of patches, jointly masking on average roughly half of
the $C \cdot P$ cells. \emph{Crucially, every $(c, p)$ cell that is
masked is masked identically across all $T$ windows}: this prevents
the predictor from exploiting the high temporal autocorrelation of
EEG to trivially copy the same channel/patch from a neighbouring
window. We reject mask draws that leave fewer than $15\%$ of tokens
visible. The context set $\mathcal{C}$ is the union of unmasked
positions; each of the $B$ blocks defines a separate prediction target
$\mathcal{M}_b$, disjoint from $\mathcal{C}$.

\paragraph{Predictor with mask tokens.}
The predictor $g_\phi$ takes the context tokens
$\{f_\theta(X)_i\}_{i \in \mathcal{C}}$, adds their position
embeddings, appends a learned mask token $\mathbf{m} \in
\mathbb{R}^{D'}$ summed with each target position embedding
$\{\mathbf{p}_j\}_{j \in \mathcal{M}_b}$, runs a transformer of the
same depth as the encoder ($L_g = L = 2$), and reads out the
final-layer states at the mask-token slots as predictions
$\hat{\mathbf{z}}_j \in \mathbb{R}^{D}$. Following I-JEPA / V-JEPA
\citep{assran2023ijepa,bardes2024vjepa}, $g_\phi$ operates at a
narrow bottleneck $D' = 24 < D = 64$ via linear input/position/output
projections, which forces the encoder to produce richer
representations rather than offloading work to the predictor.

\paragraph{Pretraining loss.}
For each masked block we minimize a smooth-L1 (Huber) regression to
the \emph{stop-gradient} target representations,
\begin{equation}
  \mathcal{L}_{\text{pred}}(X)
   = \frac{1}{|\mathcal{M}|}
     \sum_{b=1}^{B} \sum_{j \in \mathcal{M}_b}
     \mathrm{Huber}\!\Big(\hat{\mathbf{z}}_j,\;
                          \mathrm{sg}\big[f_{\bar\theta}(X)_j\big]\Big),
  \label{eq:pred-loss}
\end{equation}
with $|\mathcal{M}| = \sum_b |\mathcal{M}_b|$.

\paragraph{Anti-collapse regularization.}
A trivial constant solution drives Eq.~\eqref{eq:pred-loss} to zero,
so two complementary mechanisms prevent representational collapse.
\emph{(i) EMA target encoder.} The parameters of $f_{\bar\theta}$ are
updated as $\bar\theta \leftarrow \mu_t \, \bar\theta + (1 - \mu_t)\,
\theta$ with a cosine momentum schedule $\mu_t : 0.996 \to 1.0$ over
the full training run, and gradients flow only through $f_\theta$ and
$g_\phi$. \emph{(ii) VICReg regularization.} We additionally apply
the variance and covariance terms of VICReg \citep{bardes2022vicreg}
to the context-encoder outputs passed through a small projector
$h_\psi$,
\begin{equation}
  \mathcal{L}_{\text{vc}}(\mathbf{Z})
   = \lambda_{\sigma}\, \frac{1}{D_h}\sum_{d=1}^{D_h}
     \max\!\big(0,\, 1 - \sigma_d(\mathbf{Z})\big)
   + \lambda_{c}\, \frac{1}{D_h}\sum_{d \neq d'}
     \mathrm{Cov}_{d d'}(\mathbf{Z})^2,
   \qquad
   \mathbf{Z} = h_\psi\!\big(f_\theta(X)_{\mathcal{C}}\big),
\end{equation}
where $\sigma_d$ and $\mathrm{Cov}$ are computed across the batch and
token axes within $\mathcal{C}$ and $D_h$ is the projector output
dimension. The invariance term of VICReg is supplied implicitly by
the JEPA prediction loss in Eq.~\eqref{eq:pred-loss} and is therefore
omitted. We set $\lambda_\sigma = \lambda_c = 0.25$ and minimize
$\mathcal{L} = \mathcal{L}_{\text{pred}} + \mathcal{L}_{\text{vc}}$.

\paragraph{Optimization.}
We train for $100$ epochs with batch size $64$ on a single GPU using
AdamW \citep{loshchilov2019adamw} with peak learning rate
$5 \times 10^{-4}$, $5$-epoch linear warm-up, and cosine decay to
$10^{-6}$. Gradients on $f_\theta$, $g_\phi$, and $h_\psi$ are
clipped to global $\ell_2$ norm $1.0$; $f_{\bar\theta}$ is updated by
EMA only. At evaluation time the predictor, target encoder, and
VICReg projector are discarded. We train five seeds
$\{42, 123, 456, 789, 2025\}$ and report inter-seed variability as
our primary measure of encoder-training variance.

% =====================================================================
\subsection{Evaluation protocol}
\label{sec:method:eval}

Every method, including all baselines, is evaluated under a single
shared protocol: a deterministic clip-to-feature map followed by
closed-form linear probes and recording-level bootstrap CIs. Only
the clip-to-feature map varies between methods.

\paragraph{Per-clip representation.}
Once pretraining is complete, the context encoder $f_\theta$ is
frozen. Given an $8$\,s clip $X$ in canonical CorrCA form ($C{=}5$,
$T{=}2$, $W{=}800$), we obtain token outputs $f_\theta(X) \in
\mathbb{R}^{N \times D}$, reshape to the structured grid $[C, T, P,
D]$, and form a per-clip representation by mean-pooling over the
patch and window axes while \emph{keeping the channel axis}:
\begin{equation}
  \bar{\mathbf{h}}(X) = \mathrm{flatten}_c\!\Big(
      \tfrac{1}{T \cdot P}\!\sum_{t, p} f_\theta(X)_{c, t, p, :}
    \Big) \;\in\; \mathbb{R}^{C \cdot D}
  \;=\; \mathbb{R}^{320}.
\end{equation}
Keeping the channel axis preserves the spatial structure that CorrCA
already isolated, and makes the probe input dimension match the
linear baselines (\S\,below) for an apples-to-apples capacity
comparison.

\paragraph{Probe targets.}
We probe a deliberately broad hierarchy of targets, from low-level
sensory features through narrative structure to subject-level traits,
in order to characterise \emph{what} the representations encode rather
than only \emph{how well} they decode any one variable.
\begin{enumerate}
    \item \textit{Low-level visual features.} Per video frame we
          extract \texttt{luminance\_mean} (mean grayscale intensity)
          and \texttt{contrast\_rms} (RMS contrast), and average them
          within each window.
    \item \textit{Narrative event score} \texttt{narr} — a saliency
          score for narrative events extracted by the
          movie-annotation pipeline (\S\ref{sec:method:data}).
    \item \textit{Temporal position in the movie}
          \texttt{position\_in\_movie}, the normalised time
          $\tau / T_{\text{clip}} \in [0, 1]$, probing whether
          representations carry coarse temporal context independent
          of local stimulus content.
    \item \textit{Movie identity.} A $20$-bin top-$1$ / top-$5$
          classification over $20$ equal-width position bins, used as
          a coarser temporal probe.
    \item \textit{Subject traits.} \texttt{age} as a continuous
          regression target and biological \texttt{sex} as a binary
          classification target.
\end{enumerate}
For all probes, continuous targets are predicted by ridge regression
$\mathrm{Ridge}(\alpha{=}1)$ in closed form, binary targets by
$\ell_2$-regularized logistic regression $\mathrm{LogReg}(C{=}1)$
fitted with L-BFGS, and the $20$-bin movie-identity task by
multinomial $\ell_2$-regularized logistic regression. Regression
metrics are Pearson $r$ and (where reported) $R^2$; classification
metrics are balanced accuracy and ROC-AUC; movie-identity is reported
as top-$1$ and top-$5$ accuracy. All standardization, label
medians/means, PCA bases, and CorrCA filters are fit on training
recordings only.

\paragraph{Sampling and clip budget.}
Each recording is read $n_{\text{pass}} = 20$ times with independent
random clip starts, yielding $14{,}000$ train, $2{,}160$ val, and
$2{,}160$ test clips per encoder; val clips are used only for
protocol audits.

\paragraph{Baselines.}
Under the same protocol we compare to a linear TRF
\citep{crosse2016mtrf}, fully-supervised end-to-end EEGNet
\citep{lawhern2018eegnet}, BIOT \citep{yang2023biot}, and REVE
\citep{elouahidi2025reve}, and linear-probe + full-finetune
evaluations of LaBraM \citep{jiang2024labram}, NeuroLM
\citep{jiang2024neurolm}, and EEGPT \citep{wang2024eegpt}, all under
the same $8$\,s clip definition and R6 test split.

\paragraph{Implementation details.}
Cine-JEPA is implemented in PyTorch~$2$ \citep{paszke2019pytorch} with
\texttt{einops} for tensor reshaping, and we train each run on a
single NVIDIA GPU. All probes use scikit-learn
\citep{pedregosa2011sklearn}. Code, pretrained checkpoints, and the
movie-annotation pipeline used to compute the probe targets will be
released upon publication.
 from the main file.
% Citation keys are placeholders; resolve them in the bibliography.
% =====================================================================

\section{Method}
\label{sec:method}

\textbf{Cine-JEPA} is a joint-embedding predictive architecture
\citep{lecun2022jepa,assran2023ijepa,bardes2024vjepa,assran2025vjepa2}
that learns EEG representations by predicting the latents of masked
spatiotemporal blocks of cortical activity recorded during
naturalistic movie-watching, on top of a stimulus-locked
CorrCA-projected input. The EEG-side tokenizer, $4$D positional
embedding, and transformer trunk are taken from REVE
\citep{elouahidi2025reve} unchanged, so any gain we report is
attributable to the input pipeline, the self-supervised objective,
and our masking specialization rather than to backbone differences.

% =====================================================================
\subsection{Pretraining data and preprocessing}
\label{sec:method:data}

\paragraph{Healthy Brain Network EEG corpus.}
We pretrain on the movie-watching subset of the Healthy Brain Network
EEG dataset \citep{alexander2017hbn} accessed via the EEGDash data
loader \citep{tousignant2024eegdash} on OpenNeuro. Subjects are
children and adolescents who watched short clips while EEG was
recorded on a $129$-channel EGI GSN-HydroCel system. We use a single
clip throughout, \textit{The Present} ($203.3$\,s, $24$\,fps), to
guarantee that train, validation, and test recordings are time-aligned
to the same stimulus and therefore directly comparable under shared
movie-feature labels. HBN releases are split disjointly by subject:
releases R1--R4 (\texttt{ds005505}--\texttt{ds005508}, ${\sim}\,700$
subjects in total) for pretraining and probe training, release R5
(\texttt{ds005509}, $108$ subjects) for validation and model
selection, and release R6 (\texttt{ds005510}, $108$ subjects) as a
held-out test set used only for final evaluation. Splitting at the
release level guarantees subject-level disjointness by construction.
Recordings shorter than $10$\,s, or whose movie-onset/offset
annotations disagree with the clip duration by more than $1$\,s, are
discarded.

\paragraph{Resampling and bandpass.}
Raw recordings are resampled to a common $f_s = 200$\,Hz, bandpass
filtered to $[0.5, 99.5]$\,Hz, cast to \texttt{float32}, and clipped
to $\pm 15\sigma$ after temporal normalization to limit the influence
of movement and electrode artefacts.

\paragraph{Temporal normalization.}
Each recording $X^{(s)} \in \mathbb{R}^{C \times T_s}$ is Z-scored
\emph{per recording}: $\tilde{X}^{(s)}_{c,t} = (X^{(s)}_{c,t} -
\mu^{(s)}_c) / \sigma^{(s)}_c$, removing recording-level amplitude
differences (skin conductance, electrode impedance, instrumentation
drift) while leaving subject-trait signal recoverable by downstream
probes. Corpus-level (global) Z-scoring (the REVE convention,
\citep{elouahidi2025reve}) is reported as an ablation.

\paragraph{Spatial filtering (CorrCA).}
We then apply Correlated Component Analysis
\citep{dmochowski2012corrca} as a stimulus-locked spatial filter:
filters $W \in \mathbb{R}^{129 \times K}$ that maximize inter-subject
correlation of $W^{\!\top} X^{(s)}$ are fit on the R1--R4 recordings
of \textit{The Present}, the top $K = 5$ components are retained, and
the frozen filters are applied to all downstream data. The canonical
pipeline therefore feeds $C = 5$ CorrCA channels to the encoder;
$C = 129$ Raw is reported as an ablation.

\paragraph{Sampling and stimulus--EEG alignment.}
For each iteration we sample a contiguous $8$\,s clip ($T = 2$
windows of $4$\,s) uniformly over recordings and start times.
Cortical EEG responses to visual input are delayed: the earliest
striate response to luminance/contrast modulations sits at
${\sim}80$--$100$\,ms ($P1$ / VESPA peak) \citep{lalor2006vespa,
lalor2009vespa, lalor2012vespa}, mTRF analyses of natural stimuli
use lags out to $300$--$600$\,ms \citep{crosse2016mtrf,
diliberto2015lowfreq, lalor2010natural}, and ISC of free movie
viewing peaks at $100$--$300$\,ms \citep{dmochowski2012corrca,
dmochowski2014audience, cohen2017isc}. We therefore align labels
with a fixed visual-processing delay $\Delta = 300$\,ms: the
movie-feature label for an EEG window starting at sample $s$ is
read from frame
$k(s) = \lfloor (s / f_s - \Delta) \cdot \mathrm{fps} \rfloor$,
with boundary clamping at the movie endpoints. The exact value was
validated by a delay sweep on the training split
(App.~\ref{app:delay-sweep}).

% =====================================================================
\subsection{Cine-JEPA architecture and pretraining objective}
\label{sec:method:model}

\paragraph{Backbone instantiation.}
A clip is a tensor $X \in \mathbb{R}^{T \times C \times W}$ with
$T = 2$ non-overlapping $4$\,s windows ($W = 800$ samples at
$200$\,Hz) over $C$ channels — canonically $C = 5$ CorrCA components,
with $C = 129$ raw electrodes as an ablation. Each window is unfolded
into $P = 26$ overlapping temporal patches per channel (patch length
$\ell_p = 50$, overlap $\ell_o = 20$, stride $30$) and each patch is
linearly embedded to a $D = 64$ token, yielding $N = C \cdot T \cdot
P = 260$ tokens per clip in the canonical configuration. The context
and target encoders $f_\theta, f_{\bar\theta}$ share REVE's
transformer trunk with $L = 2$ blocks, $H = 4$ heads of dimension
$d_h = 16$, hidden width $D = 64$, and feed-forward expansion ratio
$2.66$; we refer the reader to \citep{elouahidi2025reve} for the full
block-level specification.

\paragraph{Channel-identity handling for CorrCA components.}
REVE's $4$D positional encoding takes a $3$D scalp coordinate
$(x, y, z)$ for each spatial channel together with the patch's global
time index $\tau$. For the Raw variant, $(x_c, y_c, z_c)$ is the
canonical GSN-HydroCel-$129$ coordinate of electrode $c$. The
canonical CorrCA components have no native position, so we anchor
each component $k$ to its peak-weight electrode,
\begin{equation}
  c^\star_k = \arg\max_{c} |W_{c, k}|,
  \qquad
  (x_k, y_k, z_k) = (x_{c^\star_k}, y_{c^\star_k}, z_{c^\star_k}).
\end{equation}
The CorrCA filters are dominated by a handful of central /
parieto-occipital electrodes, so the peak-weight anchor is a faithful
representative of each component's spatial support; it is also
parameter-free, montage-aware, and deterministic across runs.

\paragraph{Multi-block masking.}
We adopt V-JEPA-style multi-block masking \citep{bardes2024vjepa}
specialized to the EEG token grid. For each batch we draw $B = 4$ mask
blocks: $B_s = 2$ \emph{short-range} blocks and $B_\ell = 2$
\emph{long-range} blocks, where each block is a contiguous rectangle
on the $[C, P]$ grid. Block dimensions are sampled uniformly:
short-range blocks span $[8\%, 15\%]$ of channels and $[30\%, 60\%]$
of patches; long-range blocks span $[15\%, 35\%]$ of channels and
$[50\%, 100\%]$ of patches, jointly masking on average roughly half of
the $C \cdot P$ cells. \emph{Crucially, every $(c, p)$ cell that is
masked is masked identically across all $T$ windows}: this prevents
the predictor from exploiting the high temporal autocorrelation of
EEG to trivially copy the same channel/patch from a neighbouring
window. We reject mask draws that leave fewer than $15\%$ of tokens
visible. The context set $\mathcal{C}$ is the union of unmasked
positions; each of the $B$ blocks defines a separate prediction target
$\mathcal{M}_b$, disjoint from $\mathcal{C}$.

\paragraph{Predictor with mask tokens.}
The predictor $g_\phi$ takes the context tokens
$\{f_\theta(X)_i\}_{i \in \mathcal{C}}$, adds their position
embeddings, appends a learned mask token $\mathbf{m} \in
\mathbb{R}^{D'}$ summed with each target position embedding
$\{\mathbf{p}_j\}_{j \in \mathcal{M}_b}$, runs a transformer of the
same depth as the encoder ($L_g = L = 2$), and reads out the
final-layer states at the mask-token slots as predictions
$\hat{\mathbf{z}}_j \in \mathbb{R}^{D}$. Following I-JEPA / V-JEPA
\citep{assran2023ijepa,bardes2024vjepa}, $g_\phi$ operates at a
narrow bottleneck $D' = 24 < D = 64$ via linear input/position/output
projections, which forces the encoder to produce richer
representations rather than offloading work to the predictor.

\paragraph{Pretraining loss.}
For each masked block we minimize a smooth-L1 (Huber) regression to
the \emph{stop-gradient} target representations,
\begin{equation}
  \mathcal{L}_{\text{pred}}(X)
   = \frac{1}{|\mathcal{M}|}
     \sum_{b=1}^{B} \sum_{j \in \mathcal{M}_b}
     \mathrm{Huber}\!\Big(\hat{\mathbf{z}}_j,\;
                          \mathrm{sg}\big[f_{\bar\theta}(X)_j\big]\Big),
  \label{eq:pred-loss}
\end{equation}
with $|\mathcal{M}| = \sum_b |\mathcal{M}_b|$.

\paragraph{Anti-collapse regularization.}
A trivial constant solution drives Eq.~\eqref{eq:pred-loss} to zero,
so two complementary mechanisms prevent representational collapse.
\emph{(i) EMA target encoder.} The parameters of $f_{\bar\theta}$ are
updated as $\bar\theta \leftarrow \mu_t \, \bar\theta + (1 - \mu_t)\,
\theta$ with a cosine momentum schedule $\mu_t : 0.996 \to 1.0$ over
the full training run, and gradients flow only through $f_\theta$ and
$g_\phi$. \emph{(ii) VICReg regularization.} We additionally apply
the variance and covariance terms of VICReg \citep{bardes2022vicreg}
to the context-encoder outputs passed through a small projector
$h_\psi$,
\begin{equation}
  \mathcal{L}_{\text{vc}}(\mathbf{Z})
   = \lambda_{\sigma}\, \frac{1}{D_h}\sum_{d=1}^{D_h}
     \max\!\big(0,\, 1 - \sigma_d(\mathbf{Z})\big)
   + \lambda_{c}\, \frac{1}{D_h}\sum_{d \neq d'}
     \mathrm{Cov}_{d d'}(\mathbf{Z})^2,
   \qquad
   \mathbf{Z} = h_\psi\!\big(f_\theta(X)_{\mathcal{C}}\big),
\end{equation}
where $\sigma_d$ and $\mathrm{Cov}$ are computed across the batch and
token axes within $\mathcal{C}$ and $D_h$ is the projector output
dimension. The invariance term of VICReg is supplied implicitly by
the JEPA prediction loss in Eq.~\eqref{eq:pred-loss} and is therefore
omitted. We set $\lambda_\sigma = \lambda_c = 0.25$ and minimize
$\mathcal{L} = \mathcal{L}_{\text{pred}} + \mathcal{L}_{\text{vc}}$.

\paragraph{Optimization.}
We train for $100$ epochs with batch size $64$ on a single GPU using
AdamW \citep{loshchilov2019adamw} with peak learning rate
$5 \times 10^{-4}$, $5$-epoch linear warm-up, and cosine decay to
$10^{-6}$. Gradients on $f_\theta$, $g_\phi$, and $h_\psi$ are
clipped to global $\ell_2$ norm $1.0$; $f_{\bar\theta}$ is updated by
EMA only. At evaluation time the predictor, target encoder, and
VICReg projector are discarded. We train five seeds
$\{42, 123, 456, 789, 2025\}$ and report inter-seed variability as
our primary measure of encoder-training variance.

% =====================================================================
\subsection{Evaluation protocol}
\label{sec:method:eval}

Every method, including all baselines, is evaluated under a single
shared protocol: a deterministic clip-to-feature map followed by
closed-form linear probes and recording-level bootstrap CIs. Only
the clip-to-feature map varies between methods.

\paragraph{Per-clip representation.}
Once pretraining is complete, the context encoder $f_\theta$ is
frozen. Given an $8$\,s clip $X$ in canonical CorrCA form ($C{=}5$,
$T{=}2$, $W{=}800$), we obtain token outputs $f_\theta(X) \in
\mathbb{R}^{N \times D}$, reshape to the structured grid $[C, T, P,
D]$, and form a per-clip representation by mean-pooling over the
patch and window axes while \emph{keeping the channel axis}:
\begin{equation}
  \bar{\mathbf{h}}(X) = \mathrm{flatten}_c\!\Big(
      \tfrac{1}{T \cdot P}\!\sum_{t, p} f_\theta(X)_{c, t, p, :}
    \Big) \;\in\; \mathbb{R}^{C \cdot D}
  \;=\; \mathbb{R}^{320}.
\end{equation}
Keeping the channel axis preserves the spatial structure that CorrCA
already isolated, and makes the probe input dimension match the
linear baselines (\S\,below) for an apples-to-apples capacity
comparison.

\paragraph{Probe targets.}
We probe a deliberately broad hierarchy of targets, from low-level
sensory features through narrative structure to subject-level traits,
in order to characterise \emph{what} the representations encode rather
than only \emph{how well} they decode any one variable.
\begin{enumerate}
    \item \textit{Low-level visual features.} Per video frame we
          extract \texttt{luminance\_mean} (mean grayscale intensity)
          and \texttt{contrast\_rms} (RMS contrast), and average them
          within each window.
    \item \textit{Narrative event score} \texttt{narr} — a saliency
          score for narrative events extracted by the
          movie-annotation pipeline (\S\ref{sec:method:data}).
    \item \textit{Temporal position in the movie}
          \texttt{position\_in\_movie}, the normalised time
          $\tau / T_{\text{clip}} \in [0, 1]$, probing whether
          representations carry coarse temporal context independent
          of local stimulus content.
    \item \textit{Movie identity.} A $20$-bin top-$1$ / top-$5$
          classification over $20$ equal-width position bins, used as
          a coarser temporal probe.
    \item \textit{Subject traits.} \texttt{age} as a continuous
          regression target and biological \texttt{sex} as a binary
          classification target.
\end{enumerate}
For all probes, continuous targets are predicted by ridge regression
$\mathrm{Ridge}(\alpha{=}1)$ in closed form, binary targets by
$\ell_2$-regularized logistic regression $\mathrm{LogReg}(C{=}1)$
fitted with L-BFGS, and the $20$-bin movie-identity task by
multinomial $\ell_2$-regularized logistic regression. Regression
metrics are Pearson $r$ and (where reported) $R^2$; classification
metrics are balanced accuracy and ROC-AUC; movie-identity is reported
as top-$1$ and top-$5$ accuracy. All standardization, label
medians/means, PCA bases, and CorrCA filters are fit on training
recordings only.

\paragraph{Sampling and clip budget.}
Each recording is read $n_{\text{pass}} = 20$ times with independent
random clip starts, yielding $14{,}000$ train, $2{,}160$ val, and
$2{,}160$ test clips per encoder; val clips are used only for
protocol audits.

\paragraph{Baselines.}
Under the same protocol we compare to a linear TRF
\citep{crosse2016mtrf}, fully-supervised end-to-end EEGNet
\citep{lawhern2018eegnet}, BIOT \citep{yang2023biot}, and REVE
\citep{elouahidi2025reve}, and linear-probe + full-finetune
evaluations of LaBraM \citep{jiang2024labram}, NeuroLM
\citep{jiang2024neurolm}, and EEGPT \citep{wang2024eegpt}, all under
the same $8$\,s clip definition and R6 test split.

\paragraph{Implementation details.}
Cine-JEPA is implemented in PyTorch~$2$ \citep{paszke2019pytorch} with
\texttt{einops} for tensor reshaping, and we train each run on a
single NVIDIA GPU. All probes use scikit-learn
\citep{pedregosa2011sklearn}. Code, pretrained checkpoints, and the
movie-annotation pipeline used to compute the probe targets will be
released upon publication.
 from the main file.
% Citation keys are placeholders; resolve them in the bibliography.
% =====================================================================

\section{Method}
\label{sec:method}

\textbf{Cine-JEPA} is a joint-embedding predictive architecture
\citep{lecun2022jepa,assran2023ijepa,bardes2024vjepa,assran2025vjepa2}
that learns EEG representations by predicting the latents of masked
spatiotemporal blocks of cortical activity recorded during
naturalistic movie-watching, on top of a stimulus-locked
CorrCA-projected input. The EEG-side tokenizer, $4$D positional
embedding, and transformer trunk are taken from REVE
\citep{elouahidi2025reve} unchanged, so any gain we report is
attributable to the input pipeline, the self-supervised objective,
and our masking specialization rather than to backbone differences.

% =====================================================================
\subsection{Pretraining data and preprocessing}
\label{sec:method:data}

\paragraph{Healthy Brain Network EEG corpus.}
We pretrain on the movie-watching subset of the Healthy Brain Network
EEG dataset \citep{alexander2017hbn} accessed via the EEGDash data
loader \citep{tousignant2024eegdash} on OpenNeuro. Subjects are
children and adolescents who watched short clips while EEG was
recorded on a $129$-channel EGI GSN-HydroCel system. We use a single
clip throughout, \textit{The Present} ($203.3$\,s, $24$\,fps), to
guarantee that train, validation, and test recordings are time-aligned
to the same stimulus and therefore directly comparable under shared
movie-feature labels. HBN releases are split disjointly by subject:
releases R1--R4 (\texttt{ds005505}--\texttt{ds005508}, ${\sim}\,700$
subjects in total) for pretraining and probe training, release R5
(\texttt{ds005509}, $108$ subjects) for validation and model
selection, and release R6 (\texttt{ds005510}, $108$ subjects) as a
held-out test set used only for final evaluation. Splitting at the
release level guarantees subject-level disjointness by construction.
Recordings shorter than $10$\,s, or whose movie-onset/offset
annotations disagree with the clip duration by more than $1$\,s, are
discarded.

\paragraph{Resampling and bandpass.}
Raw recordings are resampled to a common $f_s = 200$\,Hz, bandpass
filtered to $[0.5, 99.5]$\,Hz, cast to \texttt{float32}, and clipped
to $\pm 15\sigma$ after temporal normalization to limit the influence
of movement and electrode artefacts.

\paragraph{Temporal normalization.}
Each recording $X^{(s)} \in \mathbb{R}^{C \times T_s}$ is Z-scored
\emph{per recording}: $\tilde{X}^{(s)}_{c,t} = (X^{(s)}_{c,t} -
\mu^{(s)}_c) / \sigma^{(s)}_c$, removing recording-level amplitude
differences (skin conductance, electrode impedance, instrumentation
drift) while leaving subject-trait signal recoverable by downstream
probes. Corpus-level (global) Z-scoring (the REVE convention,
\citep{elouahidi2025reve}) is reported as an ablation.

\paragraph{Spatial filtering (CorrCA).}
We then apply Correlated Component Analysis
\citep{dmochowski2012corrca} as a stimulus-locked spatial filter:
filters $W \in \mathbb{R}^{129 \times K}$ that maximize inter-subject
correlation of $W^{\!\top} X^{(s)}$ are fit on the R1--R4 recordings
of \textit{The Present}, the top $K = 5$ components are retained, and
the frozen filters are applied to all downstream data. The canonical
pipeline therefore feeds $C = 5$ CorrCA channels to the encoder;
$C = 129$ Raw is reported as an ablation.

\paragraph{Sampling and stimulus--EEG alignment.}
For each iteration we sample a contiguous $8$\,s clip ($T = 2$
windows of $4$\,s) uniformly over recordings and start times.
Cortical EEG responses to visual input are delayed: the earliest
striate response to luminance/contrast modulations sits at
${\sim}80$--$100$\,ms ($P1$ / VESPA peak) \citep{lalor2006vespa,
lalor2009vespa, lalor2012vespa}, mTRF analyses of natural stimuli
use lags out to $300$--$600$\,ms \citep{crosse2016mtrf,
diliberto2015lowfreq, lalor2010natural}, and ISC of free movie
viewing peaks at $100$--$300$\,ms \citep{dmochowski2012corrca,
dmochowski2014audience, cohen2017isc}. We therefore align labels
with a fixed visual-processing delay $\Delta = 300$\,ms: the
movie-feature label for an EEG window starting at sample $s$ is
read from frame
$k(s) = \lfloor (s / f_s - \Delta) \cdot \mathrm{fps} \rfloor$,
with boundary clamping at the movie endpoints. The exact value was
validated by a delay sweep on the training split
(App.~\ref{app:delay-sweep}).

% =====================================================================
\subsection{Cine-JEPA architecture and pretraining objective}
\label{sec:method:model}

\paragraph{Backbone instantiation.}
A clip is a tensor $X \in \mathbb{R}^{T \times C \times W}$ with
$T = 2$ non-overlapping $4$\,s windows ($W = 800$ samples at
$200$\,Hz) over $C$ channels — canonically $C = 5$ CorrCA components,
with $C = 129$ raw electrodes as an ablation. Each window is unfolded
into $P = 26$ overlapping temporal patches per channel (patch length
$\ell_p = 50$, overlap $\ell_o = 20$, stride $30$) and each patch is
linearly embedded to a $D = 64$ token, yielding $N = C \cdot T \cdot
P = 260$ tokens per clip in the canonical configuration. The context
and target encoders $f_\theta, f_{\bar\theta}$ share REVE's
transformer trunk with $L = 2$ blocks, $H = 4$ heads of dimension
$d_h = 16$, hidden width $D = 64$, and feed-forward expansion ratio
$2.66$; we refer the reader to \citep{elouahidi2025reve} for the full
block-level specification.

\paragraph{Channel-identity handling for CorrCA components.}
REVE's $4$D positional encoding takes a $3$D scalp coordinate
$(x, y, z)$ for each spatial channel together with the patch's global
time index $\tau$. For the Raw variant, $(x_c, y_c, z_c)$ is the
canonical GSN-HydroCel-$129$ coordinate of electrode $c$. The
canonical CorrCA components have no native position, so we anchor
each component $k$ to its peak-weight electrode,
\begin{equation}
  c^\star_k = \arg\max_{c} |W_{c, k}|,
  \qquad
  (x_k, y_k, z_k) = (x_{c^\star_k}, y_{c^\star_k}, z_{c^\star_k}).
\end{equation}
The CorrCA filters are dominated by a handful of central /
parieto-occipital electrodes, so the peak-weight anchor is a faithful
representative of each component's spatial support; it is also
parameter-free, montage-aware, and deterministic across runs.

\paragraph{Multi-block masking.}
We adopt V-JEPA-style multi-block masking \citep{bardes2024vjepa}
specialized to the EEG token grid. For each batch we draw $B = 4$ mask
blocks: $B_s = 2$ \emph{short-range} blocks and $B_\ell = 2$
\emph{long-range} blocks, where each block is a contiguous rectangle
on the $[C, P]$ grid. Block dimensions are sampled uniformly:
short-range blocks span $[8\%, 15\%]$ of channels and $[30\%, 60\%]$
of patches; long-range blocks span $[15\%, 35\%]$ of channels and
$[50\%, 100\%]$ of patches, jointly masking on average roughly half of
the $C \cdot P$ cells. \emph{Crucially, every $(c, p)$ cell that is
masked is masked identically across all $T$ windows}: this prevents
the predictor from exploiting the high temporal autocorrelation of
EEG to trivially copy the same channel/patch from a neighbouring
window. We reject mask draws that leave fewer than $15\%$ of tokens
visible. The context set $\mathcal{C}$ is the union of unmasked
positions; each of the $B$ blocks defines a separate prediction target
$\mathcal{M}_b$, disjoint from $\mathcal{C}$.

\paragraph{Predictor with mask tokens.}
The predictor $g_\phi$ takes the context tokens
$\{f_\theta(X)_i\}_{i \in \mathcal{C}}$, adds their position
embeddings, appends a learned mask token $\mathbf{m} \in
\mathbb{R}^{D'}$ summed with each target position embedding
$\{\mathbf{p}_j\}_{j \in \mathcal{M}_b}$, runs a transformer of the
same depth as the encoder ($L_g = L = 2$), and reads out the
final-layer states at the mask-token slots as predictions
$\hat{\mathbf{z}}_j \in \mathbb{R}^{D}$. Following I-JEPA / V-JEPA
\citep{assran2023ijepa,bardes2024vjepa}, $g_\phi$ operates at a
narrow bottleneck $D' = 24 < D = 64$ via linear input/position/output
projections, which forces the encoder to produce richer
representations rather than offloading work to the predictor.

\paragraph{Pretraining loss.}
For each masked block we minimize a smooth-L1 (Huber) regression to
the \emph{stop-gradient} target representations,
\begin{equation}
  \mathcal{L}_{\text{pred}}(X)
   = \frac{1}{|\mathcal{M}|}
     \sum_{b=1}^{B} \sum_{j \in \mathcal{M}_b}
     \mathrm{Huber}\!\Big(\hat{\mathbf{z}}_j,\;
                          \mathrm{sg}\big[f_{\bar\theta}(X)_j\big]\Big),
  \label{eq:pred-loss}
\end{equation}
with $|\mathcal{M}| = \sum_b |\mathcal{M}_b|$.

\paragraph{Anti-collapse regularization.}
A trivial constant solution drives Eq.~\eqref{eq:pred-loss} to zero,
so two complementary mechanisms prevent representational collapse.
\emph{(i) EMA target encoder.} The parameters of $f_{\bar\theta}$ are
updated as $\bar\theta \leftarrow \mu_t \, \bar\theta + (1 - \mu_t)\,
\theta$ with a cosine momentum schedule $\mu_t : 0.996 \to 1.0$ over
the full training run, and gradients flow only through $f_\theta$ and
$g_\phi$. \emph{(ii) VICReg regularization.} We additionally apply
the variance and covariance terms of VICReg \citep{bardes2022vicreg}
to the context-encoder outputs passed through a small projector
$h_\psi$,
\begin{equation}
  \mathcal{L}_{\text{vc}}(\mathbf{Z})
   = \lambda_{\sigma}\, \frac{1}{D_h}\sum_{d=1}^{D_h}
     \max\!\big(0,\, 1 - \sigma_d(\mathbf{Z})\big)
   + \lambda_{c}\, \frac{1}{D_h}\sum_{d \neq d'}
     \mathrm{Cov}_{d d'}(\mathbf{Z})^2,
   \qquad
   \mathbf{Z} = h_\psi\!\big(f_\theta(X)_{\mathcal{C}}\big),
\end{equation}
where $\sigma_d$ and $\mathrm{Cov}$ are computed across the batch and
token axes within $\mathcal{C}$ and $D_h$ is the projector output
dimension. The invariance term of VICReg is supplied implicitly by
the JEPA prediction loss in Eq.~\eqref{eq:pred-loss} and is therefore
omitted. We set $\lambda_\sigma = \lambda_c = 0.25$ and minimize
$\mathcal{L} = \mathcal{L}_{\text{pred}} + \mathcal{L}_{\text{vc}}$.

\paragraph{Optimization.}
We train for $100$ epochs with batch size $64$ on a single GPU using
AdamW \citep{loshchilov2019adamw} with peak learning rate
$5 \times 10^{-4}$, $5$-epoch linear warm-up, and cosine decay to
$10^{-6}$. Gradients on $f_\theta$, $g_\phi$, and $h_\psi$ are
clipped to global $\ell_2$ norm $1.0$; $f_{\bar\theta}$ is updated by
EMA only. At evaluation time the predictor, target encoder, and
VICReg projector are discarded. We train five seeds
$\{42, 123, 456, 789, 2025\}$ and report inter-seed variability as
our primary measure of encoder-training variance.

% =====================================================================
\subsection{Evaluation protocol}
\label{sec:method:eval}

Every method, including all baselines, is evaluated under a single
shared protocol: a deterministic clip-to-feature map followed by
closed-form linear probes and recording-level bootstrap CIs. Only
the clip-to-feature map varies between methods.

\paragraph{Per-clip representation.}
Once pretraining is complete, the context encoder $f_\theta$ is
frozen. Given an $8$\,s clip $X$ in canonical CorrCA form ($C{=}5$,
$T{=}2$, $W{=}800$), we obtain token outputs $f_\theta(X) \in
\mathbb{R}^{N \times D}$, reshape to the structured grid $[C, T, P,
D]$, and form a per-clip representation by mean-pooling over the
patch and window axes while \emph{keeping the channel axis}:
\begin{equation}
  \bar{\mathbf{h}}(X) = \mathrm{flatten}_c\!\Big(
      \tfrac{1}{T \cdot P}\!\sum_{t, p} f_\theta(X)_{c, t, p, :}
    \Big) \;\in\; \mathbb{R}^{C \cdot D}
  \;=\; \mathbb{R}^{320}.
\end{equation}
Keeping the channel axis preserves the spatial structure that CorrCA
already isolated, and makes the probe input dimension match the
linear baselines (\S\,below) for an apples-to-apples capacity
comparison.

\paragraph{Probe targets.}
We probe a deliberately broad hierarchy of targets, from low-level
sensory features through narrative structure to subject-level traits,
in order to characterise \emph{what} the representations encode rather
than only \emph{how well} they decode any one variable.
\begin{enumerate}
    \item \textit{Low-level visual features.} Per video frame we
          extract \texttt{luminance\_mean} (mean grayscale intensity)
          and \texttt{contrast\_rms} (RMS contrast), and average them
          within each window.
    \item \textit{Narrative event score} \texttt{narr} — a saliency
          score for narrative events extracted by the
          movie-annotation pipeline (\S\ref{sec:method:data}).
    \item \textit{Temporal position in the movie}
          \texttt{position\_in\_movie}, the normalised time
          $\tau / T_{\text{clip}} \in [0, 1]$, probing whether
          representations carry coarse temporal context independent
          of local stimulus content.
    \item \textit{Movie identity.} A $20$-bin top-$1$ / top-$5$
          classification over $20$ equal-width position bins, used as
          a coarser temporal probe.
    \item \textit{Subject traits.} \texttt{age} as a continuous
          regression target and biological \texttt{sex} as a binary
          classification target.
\end{enumerate}
For all probes, continuous targets are predicted by ridge regression
$\mathrm{Ridge}(\alpha{=}1)$ in closed form, binary targets by
$\ell_2$-regularized logistic regression $\mathrm{LogReg}(C{=}1)$
fitted with L-BFGS, and the $20$-bin movie-identity task by
multinomial $\ell_2$-regularized logistic regression. Regression
metrics are Pearson $r$ and (where reported) $R^2$; classification
metrics are balanced accuracy and ROC-AUC; movie-identity is reported
as top-$1$ and top-$5$ accuracy. All standardization, label
medians/means, PCA bases, and CorrCA filters are fit on training
recordings only.

\paragraph{Sampling and clip budget.}
Each recording is read $n_{\text{pass}} = 20$ times with independent
random clip starts, yielding $14{,}000$ train, $2{,}160$ val, and
$2{,}160$ test clips per encoder; val clips are used only for
protocol audits.

\paragraph{Baselines.}
Under the same protocol we compare to a linear TRF
\citep{crosse2016mtrf}, fully-supervised end-to-end EEGNet
\citep{lawhern2018eegnet}, BIOT \citep{yang2023biot}, and REVE
\citep{elouahidi2025reve}, and linear-probe + full-finetune
evaluations of LaBraM \citep{jiang2024labram}, NeuroLM
\citep{jiang2024neurolm}, and EEGPT \citep{wang2024eegpt}, all under
the same $8$\,s clip definition and R6 test split.

\paragraph{Implementation details.}
Cine-JEPA is implemented in PyTorch~$2$ \citep{paszke2019pytorch} with
\texttt{einops} for tensor reshaping, and we train each run on a
single NVIDIA GPU. All probes use scikit-learn
\citep{pedregosa2011sklearn}. Code, pretrained checkpoints, and the
movie-annotation pipeline used to compute the probe targets will be
released upon publication.
